\documentclass[12pt]{article}

\usepackage[utf8]{inputenc}
\usepackage[T1]{fontenc}
\usepackage[brazil]{babel}
\usepackage{times}
\usepackage[a4paper,left=3cm,right=3cm,top=2.8cm,bottom=2.8cm]{geometry}
\usepackage{graphicx}
\usepackage{booktabs}
\usepackage{amsmath,amssymb}
\usepackage{url}
\usepackage[numbers,sort&compress]{natbib}
\usepackage[hidelinks]{hyperref}

\newcommand{\Sset}{\mathcal{S}}
\newcommand{\Cset}{\mathcal{C}}
\newcommand{\MS}{\mathrm{EM}}
\newcommand{\taxa}{\mathrm{tx}}
\newcommand{\aprova}{\textsc{aprova}}
\newcommand{\reprova}{\textsc{reprova}}
\newcommand{\code}[1]{\texttt{\small #1}}

\title{Adequação de Suítes de Teste para Assistentes com Geração Aumentada por
Recuperação: um Estudo Baseado em Mutação}

\author{Ícaro S. de Oliveira\thanks{Universidade Estadual do Ceará (UECE),
Fortaleza, CE, Brasil. \texttt{icaro.santos@aluno.uece.br}} \and
Ismayle S. Santos\thanks{Universidade Estadual do Ceará (UECE), Fortaleza, CE,
Brasil. \texttt{ismayle.santos@uece.br}}}

\date{}

\begin{document}

\maketitle

\begin{abstract}
\noindent\textbf{Resumo.}
Assistentes inteligentes baseados em geração aumentada por recuperação (RAG)
chegaram à produção antes de existir prática consolidada para testá-los. As
suítes escritas para esses sistemas emparelham perguntas com respostas de
referência e julgam o texto gerado, enquanto os pontos de falha documentados de
RAG se concentram nas camadas de dado e de recuperação, que antecedem a geração.
Este artigo transpõe o teste de mutação para essas camadas. Derivamos 18
operadores de mutação a partir de pontos de falha documentados, cobrindo corpus,
segmentação, índice, recuperação e prompt, cada um com uma checagem executável
de fidelidade; aplicamo-los a um assistente RAG de código aberto do domínio de
engenharia de software, com uma suíte de 30 casos e quatro oráculos lendo o
mesmo log de 3.000 invocações. A suíte mata 17 dos 18 mutantes, mas detecta o
operador mediano por apenas 2,5 dos 20 casos avaliáveis, e as camadas não se
separam. O mesmo log julgado por oráculos diferentes produz escores de mutação
de 0,722 a 0,944 --- 22 pontos percentuais de oscilação --- e o oráculo mais
barato é estritamente dominado: em 32 pares (mutante, caso) as assertivas
determinísticas reprovam onde a similaridade de cosseno aprova, e em nenhum o
inverso. A campanha custou 33,9 horas, 1,08 kWh e de 63 a 83 watts-hora por
mutante morto, conforme o oráculo consultado.
\end{abstract}

\noindent\textbf{Abstract.}
Intelligent assistants built on retrieval-augmented generation reached
production before a practice of testing them matured. This paper transposes
mutation testing to the data and retrieval layers of a RAG pipeline with 18
operators derived from documented failure points, each carrying an executable
fidelity check, applied to an open-source software-engineering assistant with a
30-case suite and four oracles reading the same log of 3,000 invocations. The
suite kills 17 of the 18 mutants but detects the median operator through only
2.5 of the 20 evaluable cases. The same log judged by different oracles yields
mutation scores from 0.722 to 0.944, and the cheapest oracle is strictly
dominated.

\section{Introdução}

Assistentes inteligentes deixaram o estágio de protótipo e entraram na prática
diária da engenharia de software. Engenheiros os consultam para localizar uma
regra em uma especificação, explicar o comportamento de um framework, redigir um
caso de teste ou justificar uma decisão de projeto. A maior parte desses
assistentes é construída sobre geração aumentada por recuperação
(RAG)~\cite{lewis2020rag}: em vez de responder a partir do que o modelo
memorizou no treinamento, o sistema recupera trechos de um corpus controlado e
condiciona a resposta a eles. A arquitetura é atraente justamente porque promete
ancoragem --- respostas rastreáveis a documentos que a organização possui e
atualiza.

Essa promessa carrega uma obrigação. Um assistente que responde sobre um
contrato de interface, uma restrição regulatória ou um procedimento de teste faz
parte do processo de engenharia, e uma resposta errada não para na tela: ela se
propaga para os artefatos que o engenheiro produz a partir dela. Ainda assim,
esses sistemas chegaram à produção mais rápido do que a prática de testá-los
amadureceu.

Onde o teste existe, ele tem uma forma reconhecível. Uma suíte emparelha
perguntas com respostas de referência, executa-as contra o sistema implantado e
deriva um veredito comparando cada resposta gerada à sua referência, às vezes
reforçada por assertivas textuais ou por um segundo modelo atuando como juiz. É
a forma relatada em experiências industriais com assistentes implantados,
inclusive a nossa~\cite{oliveira2026sast}, e é a forma que um praticante adota
naturalmente, porque a resposta é o que o usuário vê.

A dificuldade é que a resposta não é onde a maioria das falhas se origina.
\citet{barnett2024seven} catalogam sete pontos de falha observados ao construir
sistemas RAG, e quatro deles ocorrem antes de um único token ser gerado:
conteúdo ausente do corpus, trecho relevante nunca recuperado, trecho recuperado
mas expulso da janela de contexto, trecho presente mas não extraído. Uma suíte
que inspeciona apenas a resposta final observa um pipeline de cinco estágios
pela sua abertura mais estreita.

Perguntar o que essa suíte deixa de ver é perguntar se ela é \emph{adequada}.
Para software determinístico o instrumento padrão é o teste de
mutação~\cite{demillo1978hints,jia2011analysis}: injetam-se defeitos pequenos no
sistema, executa-se a suíte contra cada versão defeituosa e mede-se quantos são
detectados. Transpor esse instrumento para um assistente RAG encontra dois
obstáculos.

O primeiro é que as adaptações existentes não alcançam a camada onde as falhas
estão. \citet{braga2026mutation} aplicam operadores de mutação ao prompt de
sistema de um chatbot baseado em LLM; \citet{cabral2026prompt} usam mutação para
avaliar testes unitários gerados por LLMs sob cinco estratégias de prompt;
\citet{amorim2026trigger} aplicam relações metamórficas ao prompt que dispara
agentes com chamada de ferramentas. Cada um perturba a instrução dada ao modelo
ou o código que o modelo escreve. Nenhum perturba o corpus, a estratégia de
segmentação, o índice ou os parâmetros de recuperação --- isto é, nenhum
perturba o pipeline de recuperação que distingue um sistema RAG de um chatbot.

O segundo obstáculo é que o escore de mutação repousa sobre uma suposição que
RAG viola. Diz-se que um mutante foi morto quando a suíte falha sobre ele, e o
escore só tem significado se a falha for atribuível ao defeito injetado. Um
assistente RAG devolve respostas diferentes para a mesma pergunta em execuções
distintas, de modo que um teste pode falhar por razões que nada têm a ver com a
mutação. O remédio clássico --- identificar mutantes equivalentes verificando se
o mutante se comporta exatamente como o original --- não está disponível aqui,
porque o original é ele próprio instável~\cite{atil2025nondeterminism}.

Um terceiro obstáculo apareceu enquanto construíamos os instrumentos deste
estudo, e ele muda o que um escore de adequação pode significar. O veredito que
decide se um mutante foi morto é produzido por um oráculo, e oráculos para
respostas em linguagem natural discordam entre si. Ao calibrar o mais barato e
mais difundido deles --- similaridade de cosseno entre a resposta gerada e sua
referência --- encontramos que, no nosso corpus, \emph{nenhum} limiar separa uma
paráfrase correta de um erro factual injetado (AUC $= 0{,}485$, abaixo do
acaso). O efeito não é ruído, é inversão: uma sentença com o sentido invertido
pontuou 0,9912 contra a referência, enquanto uma paráfrase fiel da mesma
referência pontuou 0,9907 e uma versão diferindo em um único valor numérico
pontuou 0,9184. A razão é estrutural: o erro injetado é a referência com um
token trocado, textualmente quase idêntica, ao passo que a paráfrase é uma
reescrita honesta.

Isso não é um problema incidental de medição. Significa que um escore de
adequação relatado sem nomear o oráculo que o produziu não é interpretável, do
mesmo modo que um escore de eficácia de reparo depende do validador que julgou
os \emph{patches}~\cite{bonfim2026evaluator}. Este estudo, portanto, relata
adequação por oráculo, e não como número único.

O sistema sob teste é um assistente RAG de código aberto cujo próprio domínio é
o teste de sistemas baseados em LLM e RAG. A escolha é deliberada: ele pertence
ao ciclo de vida da engenharia de software e é público na íntegra, corpus
incluído, de modo que o estudo possa ser replicado em vez de apenas descrito.
Fazemos três perguntas: quais camadas do pipeline produzem defeitos que a suíte
deixa de detectar (QP1), quanto a adequação relatada depende do oráculo aplicado
a artefatos idênticos (QP2) e quanto custa cada nível de rigor, em tempo de
parede, tokens, GPU-segundo e energia (QP3).

As contribuições são: (i) um catálogo de 18 operadores de mutação para pipelines
RAG, cada um derivado de um ponto de falha documentado e acompanhado de uma
checagem executável que falha no sistema não-mutado e passa no mutante; (ii) um
protocolo de adequação sob não-determinismo que separa o efeito do defeito
injetado da instabilidade intrínseca do sistema, restringindo todo denominador
aos casos estáveis no baseline; (iii) uma quantificação do viés de oráculo, com
adequação relatada por oráculo e duas das distorções estabelecidas antes da
campanha; (iv) uma caracterização de custo por nível de rigor, normalizada por
mutante morto; e (v) um pacote de replicação que roda em hardware de consumo com
modelos de pesos abertos.

\section{Fundamentação e Trabalhos Relacionados}
\label{sec:fundamentacao}

\subsection{Geração aumentada por recuperação e seus pontos de falha}

A geração aumentada por recuperação~\cite{lewis2020rag} separa o que um sistema
sabe do que seu modelo memorizou. Um corpus de documentos é segmentado em
trechos, cada trecho é projetado em um espaço vetorial e, no momento da
consulta, os trechos mais próximos da pergunta são colocados no prompt que
condiciona a geração. O pipeline tem, portanto, cinco estágios distinguíveis ---
ingestão do corpus, segmentação, indexação, recuperação e geração --- e um
defeito pode ser introduzido em qualquer um deles sem produzir erro em nenhum
outro.

\citet{barnett2024seven} relatam sete pontos de falha observados na construção
de sistemas RAG. Quatro surgem antes de o modelo escrever qualquer coisa: a
resposta está simplesmente ausente do corpus (FP1); o trecho que a contém existe
mas não é ranqueado alto o bastante para ser recuperado (FP2); é recuperado mas
não sobrevive à consolidação na janela de contexto (FP3); ou chega ao contexto e
não é extraído dele (FP4). Os três restantes dizem respeito ao texto gerado:
formato errado (FP5), nível de especificidade inadequado (FP6) e resposta
incompleta (FP7).

Essa distribuição tem uma consequência prática que motiva este estudo. Uma suíte
para assistente RAG costuma ser escrita como um conjunto de perguntas
emparelhadas com respostas de referência, e seu veredito é derivado do texto
gerado. Tal suíte observa FP5 a FP7 diretamente, porque são propriedades desse
texto. Observa FP1 a FP4 apenas quando eles por acaso mudam a resposta de um
jeito que o oráculo nota --- e uma falha de recuperação muda a resposta
sutilmente, removendo um qualificador ou um número, em vez de visivelmente.
Relatos industriais de teste de sistemas intensivos em dados descrevem suítes
exatamente com essa forma~\cite{cavalcante2025bigdata,oliveira2026sast}, o que
sugere que a assimetria não é artefato de um projeto.

A qualidade de sistemas intensivos em dados tem sido estudada de forma
sistemática na comunidade brasileira~\cite{oliveira2024bigdata}, e o padrão que
essa literatura descreve --- defeitos que nascem no dado e só aparecem no
comportamento --- é o mesmo que a transposição do teste de mutação para RAG
precisa capturar.

\subsection{O problema do oráculo em software baseado em LLM}

Decidir se uma saída está correta é o problema do oráculo, e respostas em
linguagem natural o tornam agudo: não há uma única cadeia de caracteres correta,
muitas formulações são igualmente válidas e a mesma entrada não produz a mesma
saída duas vezes. Esse último ponto vale inclusive sob configurações
nominalmente determinísticas: \citet{atil2025nondeterminism} mostram que
temperatura zero e semente fixa não eliminam variação entre execuções, o que
significa que um teste que falha não é, por si só, evidência de defeito.

\citet{dobslaw2026challenges} organizam os desafios de testar software baseado
em LLM em uma taxonomia facetada, e levantamentos da área~\cite{wang2024testing}
mapeiam as técnicas empregadas. Duas lacunas dessa literatura afetam diretamente
este estudo. A primeira é de granularidade: a taxonomia trata o sistema baseado
em LLM como a unidade sob teste, sem subdividir a camada de recuperação, de modo
que as dimensões que um praticante de RAG de fato configura --- versão do
corpus, tamanho do trecho, modelo de projeção, top-$k$, piso de relevância ---
não aparecem como facetas a controlar. A segunda é de custo: enumera-se o que
pode ser testado e como, mas não o que o teste custa em tempo, dinheiro ou
energia, embora a execução repetida seja a única resposta disponível ao
não-determinismo e repetição seja precisamente o que custa.

Na prática, quatro famílias de oráculo são usadas. A similaridade de cosseno
entre a projeção da resposta e a de uma resposta de referência é a mais barata e
a mais comum. Assertivas determinísticas sobre o texto da resposta --- presença
ou ausência de um termo, uma expressão regular, um marcador de citação --- são
ainda mais baratas, mas cobrem apenas o que pode ser enunciado como propriedade
de superfície. Pacotes de métricas sem referência, como o
RAGAS~\cite{es2024ragas}, decompõem a resposta em afirmações e checam cada uma
contra o contexto recuperado. Por fim, um segundo modelo pode atuar como
juiz~\cite{zheng2023judging}, ao custo de importar a inconsistência desse modelo
e sua preferência documentada por texto produzido por modelos da própria
família~\cite{panickssery2024llm}.

Cada uma dessas famílias carrega suposições raramente declaradas. O oráculo de
similaridade supõe que equivalência semântica é bem aproximada por proximidade
no espaço de projeção, suposição que evidência adjacente já coloca em dúvida:
representações baseadas em LLM capturam regularidades de superfície de código
muito melhor do que capturam sua semântica~\cite{infsof2025codesemantics}. Os
pacotes de métricas supõem que sua noção de fidelidade coincide com a noção de
correção que o testador tem em mente. A Seção~\ref{sec:oraculos} relata o que
acontece quando essas suposições são examinadas em vez de herdadas.

\subsection{Teste de mutação para sistemas baseados em LLM}

O teste de mutação responde a uma pergunta que a cobertura não responde: dada
uma suíte que passa, ela notaria se o sistema estivesse errado? Defeitos pequenos
são injetados, a suíte é executada contra cada versão defeituosa e a proporção
detectada --- o escore de mutação --- estima a capacidade de detecção da
suíte~\cite{demillo1978hints}. A técnica repousa sobre duas
suposições~\cite{jia2011analysis}: que um mutante morto foi morto pelo defeito
injetado e que mutantes que nenhuma suíte consegue matar são equivalentes ao
original e podem ser identificados comparando comportamentos. Ambas são estáveis
em software determinístico e nenhuma sobrevive ao contato com um assistente RAG.

Três estudos recentes transpõem injeção de defeitos para sistemas baseados em
LLM e, juntos, delimitam o que falta. \citet{braga2026mutation} aplicam
operadores ao prompt de sistema de um chatbot e relatam a consequência de a
primeira suposição falhar: um escore bruto de 0,86 cai para 0,38 quando cada
morte é inspecionada e apenas as atribuíveis ao defeito injetado são contadas. O
sistema sob teste deles não tem camada de recuperação, e a campanha para no
primeiro caso que mata, o que oculta as demais causas de morte.
\citet{cabral2026prompt} usam mutação para avaliar testes unitários gerados por
LLMs sob cinco estratégias de prompt, repetindo cada geração cinco vezes para
absorver não-determinismo e diagnosticando sobreviventes para derivar
recomendações; ali, porém, o sistema sob teste é código Java determinístico e
apenas o gerador de testes é estocástico. \citet{amorim2026trigger} tomam outro
caminho para o mesmo obstáculo: em vez de injetar defeitos, aplicam relações
metamórficas ao prompt de agentes com chamada de ferramentas e introduzem uma
relação identidade cujo único propósito é medir quanto o agente varia sem
transformação alguma, para que essa instabilidade intrínseca possa ser
descontada antes de qualquer violação ser atribuída a uma transformação.

Os três convergem para o prompt e para o código. Nenhum perturba o corpus, a
segmentação, o índice ou os parâmetros de recuperação --- isto é, nenhum perturba
os estágios onde quatro dos sete pontos de falha documentados se localizam.

\subsection{Avaliação de assistentes inteligentes}

Assistentes construídos sobre modelos de linguagem passaram a ser estudados como
artefatos de software por direito próprio~\cite{muttillo2025modeling}, e o
trabalho sobre engenharia de requisitos com LLMs coloca correção e
confiabilidade dos artefatos gerados como preocupação
central~\cite{infsof2025reroadmap}. O que essa linha estabelece é o objeto; o
que ela deixa em aberto é como o veredito sobre esse objeto é alcançado.

Três resultados afiam a pergunta. \citet{bonfim2026evaluator} mantêm os
artefatos fixos --- os mesmos 300 \emph{patches} de reparo de programa --- e
variam apenas o validador, obtendo oscilação de cerca de vinte pontos
percentuais na eficácia relatada; a qualidade relatada de um sistema é,
portanto, em parte propriedade do instrumento que a mediu. \citet{lima2026limits}
relatam um resultado delimitador em que priorização de casos de teste baseada em
LLM não supera linhas de base lexicais leves, e tratam custo como resultado de
primeira classe, postura também adotada por ferramental de \emph{benchmark} para
estratégias de provisão de contexto~\cite{vasconcelos2026llmcontextbench}.
\citet{vergilio2026rag} mostram que ganhos de RAG medidos sobre trechos curados
não se transferem para a escala de repositório, onde custo e manutenibilidade se
tornam restritivos.

Duas observações da prática industrial completam o quadro. \citet{franca2026genai}
relatam que um assistente baseado em RAG identificou sete cenários de teste para
um fluxo em que analistas humanos identificaram quarenta e oito, e atribuem a
diferença a conhecimento que existe em conversas e nunca é escrito --- um limite
que nenhum ajuste de recuperação remove. E nossa própria experiência com um
assistente implantado no setor público~\cite{oliveira2026sast} produziu uma
regressão cuja causa foi um documento normativo atualizado que permaneceu
indexado em sua versão anterior: uma falha de corpus, invisível ao monitoramento
de erros, detectada apenas porque as respostas mudaram.

\subsection{Lacuna}

Reunindo: os pontos de falha documentados de RAG concentram-se em recuperação e
dado; as suítes escritas para esses sistemas inspecionam a resposta gerada; o
teste de mutação, instrumento padrão para perguntar se uma suíte é adequada, foi
transposto para o prompt e para o código gerado, mas não para o pipeline de
recuperação; a suposição de que uma morte é atribuível ao defeito injetado não
se sustenta sob não-determinismo, e a rota clássica para identificar mutantes
equivalentes está indisponível porque o sistema não-mutado é ele próprio
instável; e o veredito que decide uma morte vem de um oráculo cujas suposições
são em geral herdadas em vez de testadas, ainda que trocar o oráculo seja
reconhecidamente capaz de mudar o escore relatado.

Nenhum estudo perguntou ainda quais camadas de um pipeline RAG produzem defeitos
que sua suíte de testes deixa de detectar, nem quanto da resposta depende do
oráculo usado para obtê-la.

\section{Contexto do Estudo}
\label{sec:contexto}

\subsection{O assistente sob teste}

O sistema sob teste é um assistente de código aberto cuja matéria é o teste de
sistemas baseados em LLM e RAG. Ele é a implementação pública de referência da
arquitetura de testes que relatamos para um assistente implantado no setor
público~\cite{oliveira2026sast}; o corpus e o conjunto de referência do sistema
original contêm informação operacional da organização parceira e não podem ser
divulgados, o que motivou construir um sistema equivalente que possa.

O pipeline roda inteiramente em modelos locais de pesos abertos. A geração usa
\code{qwen3:8b} e a projeção usa \code{nomic-embed-text}, ambos servidos pelo
Ollama; o armazenamento vetorial é o ChromaDB; o backend é uma aplicação FastAPI
que expõe consulta, reingestão do corpus e métricas de execução. Cada interação
é registrada com a pergunta, os trechos recuperados, a resposta e as métricas da
invocação. O hardware é uma GPU de notebook de consumo com 8~GB de memória.

Rodar localmente é escolha metodológica, não restrição de orçamento. Endpoints
comerciais são revisados sem aviso, o que torna irreproduzível, após o fato, uma
campanha de milhares de invocações; pesos abertos fixados por \emph{digest} não
têm esse problema. O custo é capacidade de modelo, e voltamos a ele ao discutir
validade externa.

\subsection{Corpus}

O corpus tem 8 documentos, 96 páginas e cerca de 164 mil tokens de literatura de
pesquisa sobre teste de software. Duas decisões de preparo determinam o que o
estudo pode observar.

Primeiro, \emph{todo} documento é re-renderizado a partir do seu texto extraído,
não apenas os que os operadores modificam. Se só os documentos mutados fossem
re-renderizados, a diferença de extração entre um original e sua contraparte
mutada entraria no resultado junto com o defeito injetado, e os dois seriam
indistinguíveis.

Segundo, a re-renderização repagina com orçamento de aproximadamente 1.800
tokens por página. O assistente segmenta trechos em fronteiras de página, e um
artigo em duas colunas rende páginas de cerca de 500 tokens; no tamanho de
trecho configurado, um único trecho engoliria a página inteira e os quatro
operadores de segmentação seriam inertes por construção --- não produziriam
mudança observável, e a suíte pareceria não detectar defeitos que o desenho
nunca permitiu que ela visse. Após a repaginação o índice tem 184 trechos para
96 páginas, ou 1,9 trecho por página.

A normalização de texto importou mais do que o esperado. As fontes Base-14 dos
PDFs de origem codificam apenas Latin-1, de modo que 1.372 símbolos foram
extraídos como pontos de interrogação, 677 deles em um único documento --- e
notação estatística é precisamente o que os casos factuais citam. Normalização
Unicode combinada com uma tabela de substituição para letras gregas e operadores
matemáticos, aplicada caractere a caractere, recuperou 93\% deles.

Três papéis foram atribuídos a documentos para que os operadores de corpus
tenham alvo: um documento \emph{primário}, denso em valores factuais, do qual as
variantes perturbadas derivam; um documento \emph{secundário} cujo terço final é
prosa e não bibliografia, de modo que o truncamento remova conteúdo sobre o qual
a suíte pergunta; e dois documentos \emph{recentes}, usados pelo operador que
deixa corpus e índice fora de sincronia. Duas variantes do primário foram
produzidas, uma representando versão superada e outra introduzindo valor
conflitante, cada uma com 40 substituições distribuídas pelas páginas;
marcadores de citação e entradas bibliográficas foram excluídos da perturbação,
porque nenhum caso de teste pergunta sobre eles.

O corpus é composto por artigos de pesquisa sobre teste de software, e alguns
deles são trabalhos relacionados deste próprio artigo. A sobreposição é
consequência do domínio do assistente e não entra no resultado: todo caso
dependente do corpus pergunta por um valor factual enunciado em um documento ---
uma contagem, um percentual, um limiar --- e nenhum pergunta por sua tese.

\subsection{Suíte de testes}

A suíte tem 30 casos derivados por particionamento em classes de equivalência
sobre o corpus (Tabela~\ref{tab:classes}). As classes não são conveniência
descritiva: existem para que cada camada do pipeline tenha ao menos um caso
capaz de expor um defeito nela. Sem essa garantia o estudo mediria a suíte em
vez do pipeline, e um mutante sobrevivente não poderia ser distinguido de uma
lacuna de cobertura que o próprio desenho criou.

\begin{table}[htb]
\centering
\caption{Classes de equivalência da suíte e o comportamento que cada uma
exercita.}
\label{tab:classes}
\begin{tabular}{@{}lcl@{}}
\toprule
Classe & Casos & Comportamento esperado \\
\midrule
Fato direto          & 6  & Recuperar um valor explícito \\
Fato distribuído     & 5  & Combinar evidência de dois ou mais trechos \\
Filtro condicional   & 4  & Responder conforme uma condição enunciada \\
Fora do corpus       & 5  & Abster-se \\
Fora de escopo       & 3  & Recusar \\
Ambíguo              & 3  & Pedir esclarecimento \\
Rastreabilidade      & 4  & Citar o documento de origem \\
\midrule
Total                & 30 & \\
\bottomrule
\end{tabular}
\end{table}

Cada resposta de referência foi conferida contra o trecho do índice do qual
deriva. Os 19 casos dependentes do corpus foram primeiro rascunhados pelo modelo
local a partir dos trechos indexados e depois verificados um a um contra a
fonte; em quatro rodadas nenhum número foi inventado, mas o modelo regularmente
anexava uma justificativa própria a um fato correto. Como a resposta de
referência é ela própria o oráculo, tal afirmação penalizaria uma resposta
correta do sistema; seis referências foram reescritas para conter apenas o
literal.

Duas verificações foram acrescentadas depois que casos iniciais reprovaram por
razões alheias ao que deviam testar. A primeira verifica, para todo caso
dependente do corpus, que a pergunta de fato recupera a \emph{página} que contém
sua evidência: num corpus em que todo documento discute teste, uma pergunta
genericamente formulada recupera o vizinho temático, e o caso reprovaria contra
o sistema não-mutado e sairia da análise em silêncio. A segunda executa o
oráculo determinístico contra as próprias respostas de referência e descarta
qualquer assertiva que sua própria referência não satisfaça --- uma assertiva
que falha no baseline remove o caso da análise com a mesma discrição.

\subsection{Instrumentação}
\label{sec:instrumentacao}

Operadores de mutação para um pipeline RAG mudam configuração, não código-fonte,
e vários dos parâmetros que eles mudam não existiam como valores configuráveis
no assistente. Acrescentá-los exigiu 264 linhas de mudança, cada uma introduzida
com um padrão que preserva o comportamento anterior: o prompt de sistema é
emitido byte a byte como antes, e a suíte de API pré-existente, de 18 testes,
passa inalterada.

O arcabouço nunca reescreve um arquivo sob a aplicação. Ele aplica um delta de
configuração em memória e o descarta quando o mutante é revertido, o que elimina
o modo de falha mais caro de uma campanha de mutação --- um mutante que
permaneceu parcialmente aplicado após a reversão, contaminando em silêncio toda
medição subsequente. A campanha também usa ambiente virtual e índice vetorial
próprios, de modo que executá-la não perturba o estado operacional do
assistente.

O baseline do estudo não é o padrão do repositório, e cada diferença existe para
que algum operador tenha o que remover: trechos de 1.200 tokens com sobreposição
de 200, porque os operadores de segmentação são definidos como deltas a partir
desses valores; reordenação por relevância marginal máxima ligada, porque um
operador a desliga; piso de relevância acima de zero, porque um operador o
remove; instrução de citação ligada, porque um operador a remove; e temperatura
fixada em 0,0 explicitamente, exceto no operador que a eleva. O arcabouço recusa
aplicar um operador cujo valor coincida com o do baseline, de modo que um
mutante equivalente por configuração não entra na campanha despercebido.

Um parâmetro do servidor de modelos precisou ser declarado antes que qualquer
medição fosse confiável, e ele não estava no protocolo. Sem uma janela de
contexto explícita, o servidor usa 4.096 tokens e trunca silenciosamente todo
prompt maior que a metade disso, preservando os quatro primeiros tokens e os
2.046 últimos --- o que descarta o prompt de sistema e, com ele, as regras de
ancoragem, abstenção e citação. O primeiro baseline completo deste estudo, com
300 invocações, rodou assim; o servidor só reportou isso no seu próprio log, e
todos os oráculos julgaram respostas fluentes produzidas por conhecimento
paramétrico como se as regras tivessem sido desobedecidas. A janela passou a ser
dimensionada pela configuração de recuperação de cada mutante,
$\mathrm{num\_ctx} = \lceil (k \cdot s \cdot 1{,}25 + 1.024 +
\mathrm{num\_predict}) / 1.024 \rceil \cdot 1.024$ para top-$k$ igual a $k$ e
trecho de tamanho $s$, com o teto de geração explícito porque o limite de prompt
do servidor é a diferença entre os dois. O arcabouço registra o tamanho do
prompt de toda invocação e transforma em erro auditável um prompt no limite.

Um parâmetro do baseline precisou ser calibrado antes de a campanha começar. O
piso de relevância da recuperação era valor configurado mas não aplicado, e o
operador que o remove só pode ter efeito se o baseline o aplicar. Escolher esse
piso expôs uma assimetria que vale enunciar, porque o mesmo raciocínio governa o
limiar calibrado na Seção~\ref{sec:oraculos}. Um piso alto demais rejeita uma
pergunta que deveria ter sido respondida, o caso então reprova contra o sistema
não-mutado e desaparece do denominador de toda questão de pesquisa sem que erro
algum seja levantado. Um piso baixo demais apenas deixa passar uma pergunta
inrespondível, o que a instrução de abstenção deve capturar e que é, portanto,
visível nos resultados. Os dois erros não custam o mesmo, de modo que o piso foi
posto com margem de 0,02 abaixo do caso legítimo mais fraco (0,428), mesmo com
um valor ligeiramente maior pontuando melhor no critério otimizado.

A calibração produziu também uma medida contra a qual os resultados devem ser
lidos: seis dos oito casos negativos --- as cinco perguntas fora do corpus mais
uma das três fora de escopo --- recuperam trechos \emph{acima} do piso, com
similaridade máxima entre 0,486 e 0,608. Num corpus topicamente homogêneo, todo
pergunta é do assunto, inclusive aquelas cuja resposta está ausente, e um piso
global de similaridade não distingue "coberto" de "não coberto".

\section{Operadores de Mutação para Pipelines RAG}
\label{sec:operadores}

\subsection{Derivação}

O catálogo foi derivado dos pontos de falha da Seção~\ref{sec:fundamentacao}, e não do que era
conveniente perturbar, sob cinco regras.

\emph{Um defeito por operador.} Cada operador injeta um único defeito plausível
em um único estágio do pipeline e é rastreado a um ponto de falha. Um operador
que mudasse duas coisas de uma vez produziria mortes não atribuíveis.

\emph{Todo defeito tem um reparo.} Dois operadores que se parecem são mantidos
separados quando um praticante os corrigiria de modo diferente. C1 remove um
documento do corpus, e o reparo é obter o documento de novo; E2 deixa o
documento no corpus e fora do índice, e o reparo é rodar a reingestão. A
distinção é exatamente a que importa a quem mantém o assistente.

\emph{Configuração, não código.} Defeitos de RAG são, em sua maioria, defeitos
de configuração e de dado: um tamanho de trecho escolhido para outro corpus, uma
reindexação que não rodou, um piso de relevância que nunca foi aplicado. Todo
operador é, portanto, um delta sobre a configuração do assistente ou sobre os
arquivos que ele ingere, nunca uma edição em seu código-fonte.

\emph{Dois operadores replicam incidentes.} C2, em que a versão superada de um
documento permanece indexada, e E2, em que o corpus mudou e o índice não, são as
duas falhas que observamos em operação com um assistente implantado no setor
público~\cite{oliveira2026sast}. A presença deles ancora o catálogo no que de
fato quebra, e não apenas no que poderia quebrar.

\emph{Um operador de controle.} P4 eleva a temperatura de geração de 0 para 0,8.
É o único operador sem ponto de falha associado: não injeta erro de conteúdo,
injeta instabilidade. Existe para checar o próprio protocolo --- se a regra do
conjunto avaliável faz o que promete, P4 deve aparecer como perda de vereditos
estáveis, e não como mortes coerentes.

\subsection{O catálogo}

A Tabela~\ref{tab:operadores} lista os 18 operadores por camada, com o defeito
que cada um injeta, o ponto de falha a que é rastreado, o delta de configuração
que o implementa e o número de casos de teste desenhados para poder matá-lo.

\begin{table}[htb]
\centering
\small
\caption{Catálogo de operadores de mutação. FP indica o ponto de falha de
\citet{barnett2024seven}. "Casos" é quantos casos da suíte foram desenhados para
poder matar o operador.}
\label{tab:operadores}
\begin{tabular}{@{}llllc@{}}
\toprule
Op. & Camada & Defeito injetado & Baseline $\rightarrow$ mutante & Casos \\
\midrule
C1 & Corpus       & Documento ausente (FP1)          & primário removido do corpus e do índice & 4 \\
C2 & Corpus       & Versão superada indexada (FP1)   & primário na variante anterior (40 trocas) & 2 \\
C3 & Corpus       & Extração perde o final (FP1)     & secundário com 70\% das páginas & 3 \\
C4 & Corpus       & Duplicata conflitante (FP4)      & variante conflitante como novo documento & 1 \\
K1 & Segmentação  & Perda de contexto local (FP3)    & trecho 1.200 $\rightarrow$ 400 tokens & 5 \\
K2 & Segmentação  & Queda de precisão (FP2)          & trecho 1.200 $\rightarrow$ 3.000 tokens & 4 \\
K3 & Segmentação  & Conceito partido (FP3)           & sobreposição 200 $\rightarrow$ 0 & 3 \\
K4 & Segmentação  & Corte no meio de regra (FP3)     & fronteiras deslocadas em $+600$ tokens & 5 \\
E1 & Índice       & Degradação de recall (FP2)       & projeção 768 $\rightarrow$ 384 dimensões & 7 \\
E2 & Índice       & Índice obsoleto (FP1)            & documentos recentes fora do índice & 3 \\
R1 & Recuperação  & Evidência não recuperada (FP2)   & top-$k$ 4 $\rightarrow$ 1 & 10 \\
R2 & Recuperação  & Diluição de contexto (FP3)       & top-$k$ 4 $\rightarrow$ 12 & 7 \\
R3 & Recuperação  & Trechos redundantes (FP2)        & reordenação por relevância desligada & 4 \\
R4 & Recuperação  & Trecho irrelevante (FP1)         & piso de relevância 0,428 $\rightarrow$ 0 & 5 \\
P1 & Prompt       & Conhecimento paramétrico (FP6)   & regra de ancoragem removida & 11 \\
P2 & Prompt       & Inventa em vez de abster (FP6)   & regra de abstenção removida & 7 \\
P3 & Prompt       & Perde rastreabilidade (FP5)      & regra de citação removida & 4 \\
P4 & Prompt       & Variância (controle)             & temperatura 0,0 $\rightarrow$ 0,8 & 5 \\
\bottomrule
\end{tabular}
\end{table}

Dois detalhes de desenho merecem nota. K4 mantém tamanho e sobreposição e move
apenas onde os cortes caem; é o controle que separa "o tamanho do trecho
importa" de "a posição do corte importa". E1 troca o modelo de projeção do
assistente, mas não o do oráculo de similaridade: se o instrumento de medida
mudasse junto com o sistema medido, o operador seria inobservável por
construção.

\subsection{Aplicação, reversão e fidelidade}

Aplicar um operador sempre parte de uma cópia limpa do corpus e das fontes do
índice, de modo que nenhum operador herde resíduo de outro. Um portão executado
antes da campanha aplica cada operador, reverte-o e compara a cópia de trabalho
com a assinatura limpa: os 18 revertem sem resíduo.

Reverter sem resíduo não mostra que o mutante \emph{é} o que o catálogo diz. Um
delta de configuração pode ser aceito pelo sistema e não mudar nada --- porque o
parâmetro não é aplicado, porque o corpus não tem conteúdo no ponto em que o
operador incide, ou porque o baseline já tinha o valor que o operador define.
Acoplamos, portanto, a cada operador uma checagem executável de fidelidade que
observa o efeito no estágio que o operador deveria alterar. O critério é o que
toda checagem de fidelidade deveria satisfazer: a sonda precisa \emph{falhar no
baseline e passar no mutante}. Exemplos do que se observa: 184 $\rightarrow$ 790
trechos em K1; 913 $\rightarrow$ 0 caracteres compartilhados na primeira
fronteira em K3; dimensão de vetor 768 $\rightarrow$ 384 em E1; zero trechos
devolvidos no baseline e quatro no mutante para a pergunta abaixo do piso em R4.

Duas sondas precisaram ser endurecidas depois da primeira versão, e a razão é a
mesma nas duas: uma checagem que passa pelo motivo errado é pior que checagem
nenhuma. A sonda de R4 usava a primeira pergunta fora do corpus da lista, cujo
melhor trecho já superava o piso; baseline e mutante devolviam quatro trechos e
a checagem passava sem observar nada. A sonda de K3 contava trechos, que não
mudam quando a sobreposição é removida --- o que muda é o conteúdo.

\subsection{Equivalência sob não-determinismo: o conjunto avaliável}

Em software determinístico, um mutante equivalente é aquele que nenhum teste
distingue do original, e ele é detectado comparando comportamentos. Um
assistente RAG não oferece comportamento fixo para comparar: o original responde
de modo diferente a cada execução, e "comporta-se como o original" não tem
testemunha única. Substituímos a noção por outra, computável.

Seja $\Cset$ a suíte e $v_o(m, c, r) \in \{\aprova, \reprova\}$ o veredito do
oráculo $o$ sobre o caso $c$ na execução $r$ contra o sistema $m$, com $m = b$
denotando o baseline não-mutado. O \emph{conjunto avaliável} do oráculo $o$ é o
conjunto de casos em que o baseline é ao mesmo tempo estável e correto ao longo
de suas $N_b$ execuções:
\begin{equation}
\label{eq:avaliavel}
\Sset_o = \bigl\{\, c \in \Cset \;:\; v_o(b, c, r) = \aprova \ \text{para todo}\
r = 1, \dots, N_b \,\bigr\}.
\end{equation}

Todo denominador do estudo é restrito a $\Sset_o$. Um caso instável no baseline
produziria mortes não atribuíveis a mutação alguma; um caso que reprova
estavelmente no baseline reprovaria em todo mutante e inflaria a taxa de mortes
sem carregar informação sobre o defeito injetado. As duas exclusões são
relatadas separadamente, porque significam coisas diferentes: a primeira mede a
instabilidade intrínseca do sistema, a segunda mede onde a suíte e o sistema
não-mutado já discordam.

O conjunto é por oráculo, e não global, porque a mesma execução pode ser estável
sob uma assertiva determinística e instável sob um juiz. É isso que permite à
QP2 comparar oráculos sobre artefatos idênticos: cada um é pontuado sobre o seu
próprio terreno estável.

O conjunto é também o que separa um mutante sobrevivente de um mutante que caso
algum poderia ter matado. Três fontes de equivalência por construção foram
encontradas e removidas antes da campanha, e cada uma generaliza além deste
estudo. \emph{Por estrutura do corpus}: C3 trunca os últimos 30\% de um
documento, e em um artigo de pesquisa o terço final é quase sempre a lista de
referências, sobre a qual nenhum caso pergunta. \emph{Por configuração do
baseline}: três operadores removem algo que o padrão do repositório nunca teve.
\emph{Por ingestão}: os quatro operadores de segmentação eram inertes sobre a
paginação original. A observabilidade de um operador de corpus depende da
estrutura interna do documento sobre o qual ele age, e essa dependência não
aparece em um catálogo; um catálogo de operadores para RAG precisa, portanto,
viajar acompanhado de um critério de adequação do corpus, sob pena de o estudo
reportar "a suíte não detecta" onde a verdade é "o desenho não permitia
detectar".

\section{Oráculos e Calibração}
\label{sec:oraculos}

Todo oráculo é função dos mesmos quatro artefatos registrados no log --- a
pergunta, sua resposta de referência, a resposta gerada e o contexto recuperado
--- e devolve \aprova{} ou \reprova. Os oráculos foram escolhidos para cobrir as
quatro famílias da Seção~\ref{sec:fundamentacao} e para serem os que um praticante de fato implanta,
na forma em que os implanta. O custo deles é deliberadamente assimétrico: os
dois baratos e sua conjunção são aplicados a toda execução, e os dois caros são
aplicados uma vez por par (mutante, caso), sobre a resposta modal das execuções.
A assimetria reduz as chamadas de juiz de 13.500 para 2.700; ela é parte do
resultado de custo, não um atalho escondido dele.

\textbf{O1 --- similaridade de cosseno à referência.} Aprova quando a
similaridade entre a projeção da resposta gerada e a da referência atinge um
limiar $\tau^\ast$. O modelo de projeção é fixo, independentemente da
configuração do sistema sob teste; se o oráculo seguisse o modelo do sistema, o
operador E1 mudaria o instrumento de medida junto com o objeto medido.

\textbf{O2 --- assertivas determinísticas.} Um conjunto de assertivas sobre o
texto da resposta, e apenas sobre ele. Os 30 casos carregam 71 assertivas de dez
tipos: termos exigidos, termos proibidos, expressão regular para condição
enunciada, limite de comprimento, presença ou ausência da frase de abstenção,
comportamentos de recusa e de esclarecimento e, na classe de rastreabilidade,
citação de documento e de página. O2 nunca faz assertiva sobre os trechos
recuperados. A restrição é deliberada: a hipótese por trás da QP1 é que uma
suíte típica testa a saída e não a recuperação, e uma suíte já instrumentada com
assertivas de recuperação responderia à QP1 por construção.

\textbf{O3 --- pacote de métricas sem referência.} Aplica as três métricas do
RAGAS que uma avaliação de RAG costuma relatar~\cite{es2024ragas} --- fidelidade,
relevância da resposta e precisão do contexto --- computadas por um juiz local
sobre a resposta e o contexto recuperado, aprovando quando as três atingem 0,7.

\textbf{O4 --- modelo como juiz.} Pergunta a um segundo modelo se a resposta
gerada é aceitável dada a pergunta e a referência. Três mitigações limitam o que
se pode atribuir ao juiz: ele é de família distinta da do gerador, contra a
preferência documentada de avaliadores por texto da própria
família~\cite{panickssery2024llm}; cada par é julgado cinco vezes e decidido por
maioria simples; e o juiz responde em um esquema JSON restrito, com veredito,
causa predominante e justificativa curta, de modo que sua saída possa ser lida
sem interpretação.

Uma instrução do prompt do juiz foi acrescentada após o piloto e é, ela própria,
um achado. O juiz não vê o documento de origem; vê a pergunta, a referência e a
resposta. Ele não tem, portanto, como saber se um detalhe ausente da referência
é verdadeiro, e só pode julgar se a resposta \emph{contradiz} a referência. A
primeira versão do prompt pedia que não houvesse "afirmação não sustentada", e o
juiz leu "não sustentada" como "ausente da referência": uma resposta correta e
mais detalhada que a referência reprovava por "acrescentar informação". O prompt
passou a declarar que a referência é o mínimo que uma resposta aceitável precisa
conter, não o máximo. Um juiz baseado em referência carrega viés estrutural
contra respostas \emph{melhores} que a referência, e esse viés só se torna
visível quando o sistema sob teste é bom.

\textbf{O5 --- conjunção.} $\mathrm{O5}(c) = \mathrm{O1}(c) \wedge \bigwedge_i
\mathrm{O2}_i(c)$. Não é "oráculos em paralelo": uma violação determinística
reprova o caso qualquer que seja o escore de similaridade, e uma falha semântica
o reprova qualquer que seja a assertiva. O5 é a forma a que um praticante chega
depois do primeiro falso positivo de um oráculo de similaridade, e é o oráculo
primário da análise da QP1.

\subsection{Calibração do limiar de similaridade}

Um limiar de similaridade escolhido por intuição é um oráculo não declarado. O
protocolo fixou, portanto, um procedimento de calibração antes de existir
qualquer dado de campanha: construir um conjunto rotulado de pares de respostas,
varrer $\tau \in [0{,}60; 0{,}95]$ em passos de 0,01 e tomar o $\tau^\ast$ que
maximiza o $F_1$ da classe de respostas corretas.

Cada par é uma resposta de referência e uma candidata. As candidatas da classe
\textsc{equiv} são paráfrases da referência produzidas por um modelo de família
distinta da do gerador e revisadas individualmente: as 115 paráfrases foram
lidas, 111 aprovadas e 4 rejeitadas, por degeneração ou por relação alterada. As
candidatas da classe \textsc{erro} são a referência com um fato corrompido por
regra programática --- um número substituído por outro valor plausível, uma
entidade nomeada trocada por outra do corpus, ou uma condição enunciada
invertida --- de modo que o que mudou seja conhecido por construção. Os oito
casos cuja referência é a frase fixa de abstenção foram excluídos: não existe
versão errada plausível de uma abstenção. O conjunto final tem 247 pares (111
\textsc{equiv}, 136 \textsc{erro}) sobre 22 casos.

A varredura devolveu $\tau^\ast = 0{,}60$, o mínimo da grade, com $F_1 = 0{,}62$
e área sob a curva ROC de 0,485. O $F_1$ máximo é obtido aceitando tudo: nenhum
limiar separa paráfrase correta de erro injetado. A média de similaridade das
duas classes é 0,970 para \textsc{equiv} e 0,967 para \textsc{erro}.

Não é defeito de medição. Controles diretos sobre uma mesma referência: idêntica
1,0000; paráfrase correta 0,9907; \emph{sentença negada} 0,9912; \emph{número
trocado} 0,9184; entidade trocada 0,8095; assunto alheio 0,5973; outro idioma
0,4165. Uma sentença com o sentido invertido pontua acima de uma paráfrase
legítima.

A AUC agregada esconde dois regimes opostos (Tabela~\ref{tab:auc}). Para
substituição de entidade o oráculo funciona. Para números trocados e condições
invertidas ele \emph{se inverte}: o erro injetado pontua mais alto que a
paráfrase correta em 75\% e 91\% das comparações, respectivamente. O padrão se
replica em um segundo modelo de projeção, de família diferente e metade da
dimensionalidade, de modo que o achado é sobre o \emph{método} --- cosseno
contra referência --- e não sobre a escolha de modelo deste estudo.

\begin{table}[htb]
\centering
\caption{Similaridade dos erros injetados por tipo, contra as 111 paráfrases
corretas. AUC abaixo de 0,5 significa que o erro é ordenado acima da paráfrase
correta.}
\label{tab:auc}
\begin{tabular}{@{}lcccc@{}}
\toprule
Classe da candidata & $n$ & Média & AUC (nomic) & AUC (MiniLM) \\
\midrule
Condição invertida        & 20  & 0,995 & 0,091 & 0,174 \\
Número trocado            & 53  & 0,988 & 0,248 & 0,140 \\
\emph{Paráfrase correta}  & \emph{111} & \emph{0,970} & --- & --- \\
Entidade substituída      & 63  & 0,942 & 0,808 & 0,752 \\
\bottomrule
\end{tabular}
\end{table}

A razão é estrutural. O erro injetado é a referência com um token trocado,
textualmente quase idêntica; a paráfrase é uma reescrita honesta, com outro
verbo e outra ordem. O cosseno mede proximidade de superfície, e por isso
prefere a corrupção mínima à reformulação legítima. O defeito que mais importa
em RAG --- um número errado numa sentença correta --- é exatamente o que o
oráculo mais barato e mais comum não vê.

O1 é mantido na campanha em $\tau^\ast = 0{,}60$, porque o estudo mede o que os
oráculos que praticantes implantam detectam, e é isso que esse oráculo detecta.
Duas limitações são declaradas: a classe \textsc{erro} é gerada
programaticamente e pode ser mais fácil de detectar que os erros que um modelo
comete naturalmente, e as paráfrases revisadas saíram conservadoras. As duas
empurram na mesma direção: $\tau^\ast$ e a AUC são limites superiores do
desempenho do oráculo, pelos dois lados.

\section{Desenho do Estudo}

O desenho foi registrado antes da execução: operadores, suíte, oráculos,
limiares, repetições, métricas e testes estatísticos foram fixados e
versionados. Seguimos a estrutura de relato para estudo de caso em engenharia de
software~\cite{runeson2009guidelines}, já que o objeto é um sistema estudado em
profundidade e não uma amostra de sistemas.

\subsection{Métricas}

Seja $V_o(m, c)$ o veredito do oráculo $o$ sobre o caso $c$ contra o mutante
$m$, agregado sobre as $N_m$ execuções. Com $\Sset_o$ como na
Equação~\eqref{eq:avaliavel}, a \emph{taxa de morte} de um mutante é a proporção
de casos avaliáveis que reprovam sobre ele,
\begin{equation}
\label{eq:taxa}
\taxa_o(m) = \frac{\bigl|\{\, c \in \Sset_o : V_o(m, c) = \reprova \,\}\bigr|}
{|\Sset_o|},
\end{equation}
relatada com intervalo de confiança de Wilson a 95\%. Um mutante é \emph{morto}
sob $o$ quando $\taxa_o(m) > 0$, e o \emph{escore de mutação} da suíte sob $o$ é
a proporção de mutantes mortos, $\MS_o = |\{ m : \taxa_o(m) > 0 \}| / |M|$, com
$M$ o conjunto dos 18 mutantes.

As duas leituras são mantidas separadas de propósito. Com um mutante por
operador, um escore de mutação por operador seria 0 ou 1 e seu intervalo de
confiança nada diria; a taxa de morte sobre casos tem denominador
($|\Sset_o|$) e intervalo interpretável, e é ela que entra na comparação entre
camadas. O escore de mutação é propriedade da suíte sob um oráculo, e é sobre
ele que a \emph{oscilação} da QP2 é calculada:
$\max_o \MS_o - \min_o \MS_o$.

\subsection{Protocolo de execução}

O baseline é executado $N_b = 10$ vezes por caso (300 invocações) e cada um dos
18 mutantes $N_m = 5$ vezes por caso (2.700 invocações), totalizando 3.000
invocações do assistente. A geração roda a temperatura 0,0 em todos os casos
exceto sob P4, e a semente da execução $r$ é o próprio $r$. A campanha não para
no primeiro caso que mata: todo caso roda contra todo mutante, porque a causa de
cada morte é um resultado e parar cedo a esconderia.

Dentro de cada repetição, a ordem dos casos é embaralhada com o índice da
repetição como semente. A razão é um achado do baseline: a temperatura 0 a
semente não muda nada, mas a resposta depende do prompt que a precedeu no slot
do servidor, por reuso de prefixo no cache de chaves e valores. Em ordem fixa,
as repetições 2 a 10 de um caso seguem sempre o mesmo antecessor e saem byte a
byte iguais, enquanto a repetição 1 difere; dez repetições amostram dois
estados. Embaralhar dá a cada invocação um antecessor diferente, que é o que
repetir deveria significar em um servidor que mantém cache.

Os oráculos nunca rodam durante a campanha. Eles leem o log depois, e cada
veredito registra qual oráculo, qual nível de rigor e de quantas execuções foi
derivado. Essa separação é o que torna gratuito recalibrar um limiar, trocar a
implementação de um oráculo ou reexecutar o juiz.

\subsection{Níveis de rigor e custo}

A QP3 compara três níveis derivados do mesmo log. L1 é execução única: o
veredito da primeira execução, sem conjunto avaliável --- é o que um protocolo
que executa cada caso uma vez teria relatado, e é o que a maioria das suítes na
prática relata. L2 é o protocolo registrado da campanha: cinco execuções por par
com veredito por maioria, e $\Sset_o$ a partir das dez execuções do baseline. L3
acrescenta 30 execuções de baseline sobre um subconjunto de dez casos, para
estimar quanto o juízo de estabilidade de L2 mudaria com seis vezes mais
repetições.

Custo é medido, não estimado. Para cada invocação, o arcabouço registra tempo de
parede, tokens de prompt e de resposta e, amostrando a GPU pela biblioteca de
gerenciamento do fabricante durante a chamada, GPU-segundo e energia em
watts-hora. Os oráculos são medidos do mesmo modo.

\subsection{Análise estatística}

Todos os testes foram fixados antes da coleta, a $\alpha = 0{,}05$; nenhuma
proporção é relatada sem intervalo de Wilson a 95\% e nenhuma diferença de média
sem intervalo por \emph{bootstrap} percentílico a 95\%. Para a QP1, as taxas de
morte por operador são agrupadas por camada e comparadas por Kruskal-Wallis, com
$\eta^2$ como tamanho de efeito e comparações de Dunn sob correção de
Holm~\cite{holm1979simple}. Para a QP2, os vereditos binários dos oráculos sobre
os mesmos pares são comparados pelo $Q$ de Cochran~\cite{cochran1950comparison},
seguido de testes de McNemar par a par sob Holm, e a concordância é resumida por
$\kappa$ de Fleiss e por $\kappa$ de Cohen par a par. Para a QP3, o custo por
mutante morto em cada nível é relatado com seu intervalo por \emph{bootstrap}.

\subsection{Desvios do protocolo registrado}

Cinco desvios ocorreram, todos antes de qualquer resultado de mutante ser
examinado.

\emph{Janela de contexto.} O primeiro baseline completo rodou com o prompt
truncado, como descrito na Seção~\ref{sec:instrumentacao}. As 300 invocações foram arquivadas, não
descartadas, e repetidas com a janela dimensionada por configuração.

\emph{Checagem de recuperação em nível de página.} Um diagnóstico que não depende
de geração mostrou que em dez dos 19 casos factuais o documento certo era
recuperado, mas não a página que contém a evidência, e que isso é propriedade da
pergunta e não da configuração. A checagem foi estendida de documento para
página, e oito perguntas foram reescritas com termos do trecho da evidência.
Três casos não têm formulação natural que vença as páginas vizinhas do mesmo
documento; permanecem como estão, marcados como lacuna aceita, e saem de
$\Sset_o$ pela Equação~\eqref{eq:avaliavel}.

\emph{Instrução de citação.} Solicitado a citar "no formato [documento, pág.
N]", o gerador copiava a palavra \emph{documento} literalmente. Uma regra que o
baseline não consegue cumprir torna equivalente por construção o operador que a
remove, de modo que a regra passou a nomear o arquivo e a dar um exemplo.

\emph{Oráculo O3.} Nenhuma das duas implementações do pacote de métricas
produziu vereditos utilizáveis com o juiz local. A implementação de referência
não consegue interpretar a saída do juiz em seis de nove avaliações. A
implementação alternativa avalia tudo, mas devolve fidelidade 0,0 para uma
resposta ancorada e 1,0 para uma resposta que inventa um número sobre um tópico
ausente do contexto recuperado. Com um juiz do dobro do tamanho, a resposta
inventada continua com fidelidade 1,0: a métrica conta afirmações que não
\emph{contradizem} o contexto, e uma afirmação sobre tópico ausente não
contradiz nada. O ponto cego é por definição, não por tamanho de modelo, e o
defeito que o operador P2 injeta --- responder em vez de abster-se --- é
invisível a ele. A contingência registrada foi aplicada: a campanha rodou com
quatro oráculos, e O3 é relatado como ponto cego medido em vez de instrumento
produtor de vereditos. Como julgar é separado de executar, a escolha é
reversível sem gastar GPU novamente.

\emph{Teto de geração.} Quarenta e uma invocações da campanha (1,5\%) gastaram
todo o orçamento de geração de 4.096 tokens no traço de raciocínio e devolveram
resposta vazia. São registradas como erro e excluídas dos vereditos ---
reprová-las atribuiria ao mutante um efeito do nosso próprio teto --- e
reexecutadas uma vez por uma passada explícita de repetição, após a qual todas
completaram. A passada de repetição é um desvio que vale nomear: ela reexecuta
apenas o que falhou, de modo que, se o estouro do teto fosse correlacionado ao
efeito de um mutante, o procedimento enviesaria contra detectá-lo. As 41
concentram-se em três casos, e um deles está fora de $\Sset_o$ para todo oráculo
de qualquer modo.

A regra de corte do cronograma --- descartar os cinco operadores marcados como
cortáveis, ou reduzir a suíte para 20 casos --- \emph{não} foi aplicada: o
desenho completo de $18 \times 30 \times 5$ rodou, e o custo é relatado como
resultado em vez de absorvido como corte.

\section{Resultados}
\label{sec:resultados}

\subsection{O conjunto avaliável}

A campanha são 3.000 invocações do assistente: 300 do baseline não-mutado e
2.700 dos mutantes, sem invocação falha no log final. A
Tabela~\ref{tab:avaliavel} reporta, por oráculo, quanto da suíte o protocolo
consegue medir.

\begin{table}[htb]
\centering
\caption{Conjunto avaliável por oráculo. Um caso sai de $\Sset_o$ ou porque seu
veredito de baseline é instável ao longo das dez execuções, ou porque reprova
estavelmente no sistema não-mutado.}
\label{tab:avaliavel}
\begin{tabular}{@{}lcccl@{}}
\toprule
Oráculo & $|\Sset_o|$ & Instáveis & Reprovam & Flakiness [IC 95\%] \\
\midrule
O1 cosseno      & 26 & 1 & 3 & 0,033 [0,006; 0,167] \\
O2 assertivas   & 20 & 3 & 7 & 0,100 [0,035; 0,256] \\
O4 juiz         & 22 & 0 & 8 & 0,000 [0,000; 0,114] \\
O5 conjunção    & 20 & 3 & 7 & 0,100 [0,035; 0,256] \\
\bottomrule
\end{tabular}
\end{table}

Três leituras importam. Primeiro, os oráculos discordam sobre o que é mensurável
antes de julgar um único mutante: o de similaridade mantém 26 dos 30 casos e o
determinístico mantém 20, porque um oráculo mais estrito reprova mais casos no
sistema não-mutado. Um estudo que relatasse adequação sem relatar $|\Sset_o|$
estaria comparando escores calculados sobre denominadores diferentes.

Segundo, a instabilidade residual do sistema sob teste é pequena mas não nula, e
depende do oráculo: um caso para O1, três para O2 e O5, nenhum para O4. O zero
do juiz é propriedade do juiz, não do assistente --- ele agrega a resposta modal
e vê um único texto por par, de modo que não observa a variação que os oráculos
por execução observam.

Terceiro, os casos que reprovam estavelmente no baseline reprovam por razões que
sabemos nomear, e nomeá-las é o que mantém a exclusão honesta: três recuperam o
documento certo mas não a página da evidência; um pergunta sobre uma tabela com
nove irmãs quase idênticas no mesmo documento; um é respondido de forma
incompleta; três fazem pergunta ambígua e recebem abstenção em vez de pedido de
esclarecimento --- o prompt do assistente não tem regra de esclarecimento --- e
um copia o exemplo dado na regra de citação em vez do nome do documento. Nenhum
é defeito que a campanha injetou.

\subsection{QP1 --- adequação por camada}

Sob o oráculo primário, a suíte mata 17 dos 18 mutantes ($\MS_{O5} = 0{,}944$,
IC 95\% [0,742; 0,990]). O único sobrevivente é R4, o operador que remove o piso
de relevância da recuperação.

O escore de mutação, porém, é o número menos informativo da
Tabela~\ref{tab:taxa}. O que a taxa de morte mostra é que a suíte detecta cada
defeito por pouquíssimos de seus casos: o operador mediano é morto por 2,5 dos
20 casos avaliáveis, e apenas três operadores --- P2, com a regra de abstenção
removida (0,450); C1, com o documento primário removido (0,400); e R1, com
top-$k$ reduzido a um (0,400) --- são detectados por mais de um terço da suíte.
Um defeito injetado no corpus ou no índice é visível a uma suíte que julga
respostas apenas pelos casos cuja evidência ele por acaso toca.

\begin{table}[htb]
\centering
\caption{Taxa de morte por operador sob O5, sobre os 20 casos avaliáveis, com
intervalo de Wilson a 95\%.}
\label{tab:taxa}
\begin{tabular}{@{}llcl@{}}
\toprule
Op. & Camada & Casos & Taxa de morte [IC 95\%] \\
\midrule
P2 & Prompt       & 9 & 0,450 [0,258; 0,658] \\
C1 & Corpus       & 8 & 0,400 [0,219; 0,613] \\
R1 & Recuperação  & 8 & 0,400 [0,219; 0,613] \\
E1 & Índice       & 6 & 0,300 [0,145; 0,519] \\
K4 & Segmentação  & 5 & 0,250 [0,112; 0,469] \\
K2 & Segmentação  & 4 & 0,200 [0,081; 0,416] \\
P3 & Prompt       & 4 & 0,200 [0,081; 0,416] \\
E2 & Índice       & 3 & 0,150 [0,052; 0,360] \\
R2 & Recuperação  & 3 & 0,150 [0,052; 0,360] \\
C3 & Corpus       & 2 & 0,100 [0,028; 0,301] \\
K1 & Segmentação  & 2 & 0,100 [0,028; 0,301] \\
P1 & Prompt       & 2 & 0,100 [0,028; 0,301] \\
R3 & Recuperação  & 2 & 0,100 [0,028; 0,301] \\
C2 & Corpus       & 1 & 0,050 [0,009; 0,236] \\
C4 & Corpus       & 1 & 0,050 [0,009; 0,236] \\
K3 & Segmentação  & 1 & 0,050 [0,009; 0,236] \\
P4 & Prompt       & 1 & 0,050 [0,009; 0,236] \\
R4 & Recuperação  & 0 & 0,000 [0,000; 0,161] \\
\bottomrule
\end{tabular}
\end{table}

Agrupadas por camada, as taxas não diferem: índice 0,225; prompt 0,200;
recuperação 0,163; corpus 0,150; segmentação 0,150, com Kruskal-Wallis
$H = 1{,}368$, $p = 0{,}850$ sobre cinco grupos de no máximo quatro operadores,
e nenhuma comparação par a par sobrevivendo a Holm. A hipótese de que os
sobreviventes se concentrariam nas camadas de dado e recuperação não se
sustenta; com no máximo quatro operadores por camada o teste tem pouco poder, de
modo que a leitura honesta é que este desenho não separa as camadas, e não que
as camadas sejam equivalentes.

Dois operadores merecem relato próprio, porque em ambos o número esconde o
mecanismo.

\emph{R4 sobrevive por uma razão que o corpus explica.} Os cinco casos que o
miram são as perguntas fora do corpus, e a Seção~\ref{sec:instrumentacao} já mostrou que, num corpus
topicamente homogêneo, essas perguntas recuperam quatro trechos bem acima de
qualquer piso que preserve os casos legítimos. Remover o piso, portanto, não
muda nada: o mutante é equivalente \emph{neste corpus}, e o sobrevivente é
propriedade do corpus e da suíte, não lacuna do teste do assistente. Um estudo
que relatasse R4 como defeito não detectado estaria relatando um artefato.

\emph{E1 não degrada a recuperação, ele a interrompe.} Trocar o modelo de
projeção por um menor e de outra família deixa 95 de suas 150 invocações com
\emph{zero} trechos recuperados e prompt mediano de 302 tokens, contra 4.646 no
baseline. A causa é o piso de relevância: 0,428 foi calibrado na escala de
similaridade de um modelo, e a escala do outro modelo é diferente. O piso é
propriedade do par (corpus, modelo de projeção), não do corpus, e uma
atualização de rotina do modelo o invalida em silêncio. É o achado da
Seção~\ref{sec:instrumentacao}
visto do outro lado: lá o piso deixava de rejeitar uma pergunta que deveria
rejeitar; aqui ele rejeita a evidência que deveria manter.

\subsection{QP2 --- viés de oráculo}

Sobre os 342 pares (mutante, caso) avaliáveis nos quatro oráculos, os quatro
discordam ($Q$ de Cochran $= 58{,}62$, $p = 1{,}2 \times 10^{-12}$), com
$\kappa$ de Fleiss $= 0{,}681$ --- concordância substancial na leitura usual ---
e ainda assim uma oscilação de 22,2 pontos percentuais no escore de mutação
relatado (Tabela~\ref{tab:vies}). A mesma suíte, as mesmas 3.000 invocações, os
mesmos artefatos: 0,722 sob o oráculo mais barato e 0,944 sob o determinístico.
A oscilação é da mesma ordem dos vinte pontos relatados para validadores de
reparo de programa~\cite{bonfim2026evaluator} sobre outra tarefa, o que sugere
efeito de avaliação baseada em referência, e não deste pipeline.

\begin{table}[htb]
\centering
\caption{Adequação por oráculo sobre artefatos idênticos. Oscilação $= \max -
\min = 0{,}222$.}
\label{tab:vies}
\begin{tabular}{@{}lccl@{}}
\toprule
Oráculo & $|\Sset_o|$ & Mortos & $\MS$ [IC 95\%] \\
\midrule
O1 cosseno    & 26 & 13/18 & 0,722 [0,491; 0,875] \\
O4 juiz       & 22 & 16/18 & 0,889 [0,672; 0,969] \\
O2 assertivas & 20 & 17/18 & 0,944 [0,742; 0,990] \\
O5 conjunção  & 20 & 17/18 & 0,944 [0,742; 0,990] \\
\bottomrule
\end{tabular}
\end{table}

Os testes par a par dizem algo mais forte que a oscilação. Em 32 pares o oráculo
determinístico reprova onde o de similaridade aprova, e em \emph{zero} pares o
inverso ocorre (McNemar, $p < 10^{-4}$ após Holm): neste corpus, O1 é
estritamente dominado por O2. A conjunção O5, definida como
$\mathrm{O1} \wedge \mathrm{O2}$, é portanto idêntica a O2 em cada um dos 342
pares --- seu componente de similaridade nunca contribui um veredito. A forma
prática a que um praticante chega depois do primeiro falso positivo de um
oráculo de similaridade é, neste corpus, as assertivas sozinhas.

O juiz fica entre os dois, e não por ser versão fraca de nenhum: reprova 4 pares
que as assertivas aprovam e aprova 22 que elas reprovam ($p = 0{,}002$ após
Holm). Lendo esses 22: a assertiva exige um valor literal e o juiz aceita uma
resposta que enuncia o valor em outra forma ou omite um segundo número menos
central. Nenhum instrumento está certo em abstrato --- eles codificam noções
diferentes de "correto", e a noção é o que um estudo precisa declarar.

A discordância não é uniforme ao longo da suíte. Ela se concentra em
rastreabilidade (0,222 dos pares) e em recusa fora de escopo (0,167), e é menor
na classe fora do corpus (0,056), onde o marcador de abstenção torna o veredito
mecânico para todos os oráculos ao mesmo tempo. Prevíamos abstenção e
rastreabilidade; rastreabilidade se confirma e abstenção é a classe em que os
oráculos mais concordam, porque a frase fixa de abstenção é exatamente o tipo de
propriedade de superfície que assertiva, escore de similaridade e juiz leem do
mesmo modo.

\subsection{QP3 --- custo do rigor}

A campanha no nível de rigor registrado custou 33,9 horas de tempo de parede,
13,6 milhões de tokens, 47.326 GPU-segundos e 1,08 kWh, medidos por invocação
pelos contadores da própria GPU. A unidade é 27 segundos e 0,36 watt-hora por
invocação do assistente; o traço de raciocínio que o gerador emite antes de cada
resposta é a maior parte disso.

Normalizado pelo que comprou (Tabela~\ref{tab:custo}), um mutante morto custa
cerca de duas horas de parede e 63 watts-hora sob o oráculo determinístico, e um
terço a mais sob o de similaridade, que mata quatro mutantes a menos pela mesma
execução. O juiz acrescenta 11.109 segundos e 98 watts-hora próprios ao longo da
campanha --- 9\% sobre o custo do assistente --- pelos dois mutantes que o
oráculo de similaridade não vê.

\begin{table}[htb]
\centering
\caption{Custo da campanha e custo por mutante morto. A execução é compartilhada
por todos os oráculos; apenas o juiz acrescenta custo próprio mensurável.}
\label{tab:custo}
\begin{tabular}{@{}lcccc@{}}
\toprule
Oráculo & Parede (h) & GPU-s & Wh & Wh por morte \\
\midrule
O1 & 33,9 & 47.326 & 1.079 & 83,0 \\
O2 & 33,9 & 47.326 & 1.079 & 63,5 \\
O4 & 37,0 & 51.183 & 1.177 & 73,5 \\
O5 & 33,9 & 47.326 & 1.079 & 63,5 \\
\bottomrule
\end{tabular}
\end{table}

\section{Discussão}

\subsection{Adequação é propriedade do par (suíte, oráculo)}

A mesma suíte, as mesmas 3.000 invocações, os mesmos 342 pares: escore de
mutação de 0,722 se o veredito vem da similaridade de cosseno e de 0,944 se vem
das assertivas determinísticas. Vinte e dois pontos de adequação são propriedade
do instrumento, não do assistente, e a direção foi prevista pela calibração
antes de a campanha rodar --- O1 não consegue reprovar um número errado, e o
juiz o reprova em todo controle. Um escore de mutação para assistente RAG
relatado sem o oráculo que o produziu não é interpretável, e sugerimos que
estudos de sistemas baseados em LLM relatem adequação por oráculo como prática
corrente, como estudos de reparo de programa começaram a relatar eficácia por
validador~\cite{bonfim2026evaluator}.

\subsection{Para que serve o oráculo mais barato}

A similaridade de cosseno a uma referência não é inútil; ela está mal
especificada como oráculo de correção. Nos nossos dados ela separa tópico com
folga e é quase cega a valor e a polaridade. Detecta, portanto, as falhas que
mudam \emph{sobre o que} a resposta é --- uma abstenção onde cabia resposta, um
desvio para assunto vizinho, uma recusa --- e perde as que mudam o que a resposta
\emph{diz}. A forma prática é a conjunção: um portão de similaridade para
desvio e assertivas determinísticas para cada valor de que o caso depende. Onde
existe resposta de referência, os números nela são as assertivas.

\subsection{A resposta de referência é o oráculo}

Três das verificações acrescentadas durante este estudo existem porque uma
resposta de referência não é documentação, é instrumento. Uma referência com
justificativa que o modelo anexou a um fato correto penaliza uma resposta
correta; uma assertiva que sua própria referência não satisfaz remove o caso da
análise em silêncio; um juiz que lê a referência como teto reprova respostas
melhores que ela. As três são invisíveis até que o oráculo seja executado contra
o sistema não-mutado, e as três enviesam na mesma direção --- para reprovar o
baseline, o que, pela Equação~\eqref{eq:avaliavel}, encolhe o denominador sem
levantar erro. A prática mais útil que podemos recomendar a quem constrói uma
suíte para assistente RAG é executar todo oráculo contra o baseline antes de
confiar nele contra qualquer outra coisa.

\subsection{Operadores de mutação precisam de um critério de corpus}

Um operador de mutação para código-fonte carrega sua própria observabilidade: um
operador alterado é exercitado se a linha é executada. Um operador de mutação
para corpus não. Se truncar um documento remove algo sobre o que a suíte
pergunta, se uma mudança de tamanho de trecho move alguma fronteira, se um piso
de similaridade rejeita alguma coisa --- tudo isso depende do corpus, da
paginação e da configuração do baseline. Três dos nossos dezoito operadores eram
inertes na primeira versão do corpus e nenhum deles teria reportado erro; eles
sobreviveriam, e a sobrevivência seria lida como lacuna da suíte. Um catálogo de
operadores para RAG é, portanto, metade do instrumento; a outra metade é o
conjunto de medidas --- densidade de fatos por página, trechos por página,
distribuição das similaridades dos casos negativos --- que estabelece, antes da
campanha, que cada operador tem sobre o que agir.

\subsection{Corpus homogêneo e o piso de relevância}

Um assistente sobre corpus de um único domínio é o caso industrial comum: as
normas de um órgão, os manuais de um produto, os artigos de uma área. Num corpus
assim toda pergunta é do assunto, inclusive aquelas cuja resposta está ausente,
e um piso global de similaridade não distingue cobertura de proximidade
temática. A consequência para o teste é que a classe fora do corpus de uma suíte
exercita a instrução de abstenção do prompt muito mais do que o piso de
recuperação, de modo que uma suíte que parece testar o \emph{fallback} de
recuperação está, na verdade, testando conformidade de prompt. A consequência
para a engenharia é que abstenção precisa ser propriedade da geração, e não da
recuperação, em qualquer corpus com essa forma. E o piso é frágil nos dois
sentidos: o operador E1 mostrou que trocar o modelo de projeção o transforma em
um filtro que rejeita tudo. Ele pertence à família de propriedades do dado que
são invisíveis até que se construa um instrumento para procurá-las, e, como
elas, é barato de medir uma vez e caro de descobrir em produção.

\subsection{O parâmetro que ninguém declara}

A janela de contexto do servidor de modelos não está na lista de parâmetros a
controlar da literatura de RAG, e não estava no nosso protocolo. Ela decidiu o
resultado de 318 invocações antes que lêssemos uma linha de log. Uma suíte que
julga respostas não consegue vê-la, porque o modelo responde fluentemente a
partir do que memorizou quando as regras e a evidência são cortadas; uma
checagem de fidelidade sobre os operadores também não, porque os operadores
funcionaram --- sobre um prompt que o modelo nunca leu. O único instrumento que
a teria capturado no primeiro dia é o que acrescentamos depois: uma regra
sentinela no prompt de sistema cuja ausência na resposta significa que o prompt
não chegou. Sugerimo-lo como checagem permanente para qualquer estudo de sistema
baseado em LLM, ao lado da declaração da própria janela.

\subsection{Que camadas a suíte não vê}

A hipótese que registramos --- que os sobreviventes se concentrariam nas camadas
de dado e de recuperação --- não é o que a campanha mostra. Dezessete de dezoito
mutantes morrem, as camadas não se separam, e o único sobrevivente é explicado
pelo corpus e não pela suíte. O resultado que a substitui é sobre \emph{quão
fino} é o fio da detecção: o operador mediano é pego por 2,5 de 20 casos, e um
defeito em FP1 ou FP2 só se torna visível pelos casos cuja evidência ele por
acaso toca. C1 remove o documento primário inteiro: quatro casos foram
desenhados para pegá-lo e oito notam. C2 substitui esse mesmo documento por uma
versão superada e \emph{um} caso nota, porque só um pergunta por um valor em que
as duas versões discordam. O assistente que estudamos em
produção~\cite{oliveira2026sast} falhou exatamente do segundo modo --- documento
normativo atualizado na origem e deixado indexado na versão anterior --- e uma
suíte com esta forma tinha um caso em vinte que poderia tê-lo visto.

Esta é a leitura prática da QP1, e ela não depende da comparação entre camadas:
um escore de mutação perto de um pode conviver com uma suíte que detecta cada
defeito de camada de dado por uma única pergunta. Relatar a taxa de morte ao
lado do escore é o que torna a diferença visível, e não custa nada --- o
denominador já está ali.

\section{Ameaças à Validade}

Seguimos as quatro categorias usuais~\cite{runeson2009guidelines} e distinguimos
as ameaças que antecipamos daquelas que o estudo demonstrou; estas últimas são
enunciadas com a evidência, porque uma ameaça demonstrada é também um resultado.

\textbf{Validade de constructo.} O constructo central é "a suíte detecta o
defeito", e ele é operacionalizado por um oráculo. A Seção~\ref{sec:oraculos}
mostra que o oráculo mais barato não mede correção neste corpus, mas proximidade
temática, e que o pacote de métricas padrão não mede invenção sobre tópico
ausente. Não tratamos isso como ameaça a mitigar, e sim como achado; mitigamos a
ameaça de constructo relatando adequação por oráculo e nunca como número único.
Dois instrumentos permanecem como limites superiores: $\tau^\ast$ foi calibrado
sobre referências corrompidas programaticamente e paráfrases conservadoras, e o
juiz vê a referência e não o documento, de modo que só pode detectar contradição
e não invenção compatível com a referência. As respostas de referência são elas
próprias um constructo: foram rascunhadas por um modelo e conferidas uma a uma
contra a fonte, e seis foram reescritas para remover justificativas que o modelo
havia acrescentado.

\textbf{Validade interna.} Uma morte precisa ser atribuível ao defeito injetado.
Três mecanismos tratam disso: a checagem de fidelidade de cada operador, que
descarta um mutante que não muda nada; a checagem de resíduo, que descarta um
mutante que contamina o seguinte; e o conjunto avaliável, que descarta mortes
sobre casos que o sistema não-mutado não responde de forma estável. O juiz é de
família distinta da do gerador, contra autopreferência; o parafraseador do
conjunto de calibração é de família distinta da do gerador, contra
circularidade; ambos são da mesma família entre si, e esse é risco residual que
declaramos. O baseline truncado é a ameaça interna mais séria que o estudo
encontrou, e não foi capturada por nenhum oráculo, portão ou teste: todos os
instrumentos julgavam a resposta, e o defeito estava no que o modelo recebia.
Relatamo-lo porque uma replicação em outro servidor ou modelo encontraria o
mesmo padrão silencioso.

\textbf{Validade externa.} O estudo é de um assistente, um corpus de oito
documentos em um domínio, uma suíte de 30 casos e um par de modelos de pesos
abertos de 8 e 4 bilhões de parâmetros em hardware de consumo. O catálogo e o
protocolo pretendem transferir; os números, não. Dois dos achados foram
checados quanto à dependência do instrumento e resistiram: a inversão do oráculo
de similaridade se replica em um segundo modelo de projeção de outra família, e
o ponto cego do pacote de métricas persiste sob um juiz do dobro do tamanho. O
corpus é topicamente homogêneo, que é o caso industrial comum e também o que
derrota um piso global de relevância; um corpus heterogêneo moveria detecções da
camada de prompt para a de recuperação, e o resultado da QP1 deve ser lido com
isso em mente.

\textbf{Validade de conclusão.} Há um mutante por operador, de modo que o escore
de mutação por operador é uma única observação binária; toda afirmação
inferencial é feita sobre a taxa de morte por casos, cujo denominador é
$|\Sset_o|$ e cujo intervalo é relatado. A suíte tem 30 casos e as classes têm
entre três e seis; a comparação por camada agrupa no máximo quatro operadores, e
Kruskal-Wallis sobre grupos desse tamanho tem pouco poder, o que relatamos como
tal em vez de compensar. Um operador é mirado por um único caso, e qualquer
afirmação sobre ele repousa sobre uma observação. As cinco repetições por
mutante não são sorteios independentes da distribuição de respostas do sistema:
a temperatura 0 elas diferem apenas pelo estado do cache do servidor, que a
ordem embaralhada varia mas não aleatoriza, e a estabilidade que medem é um
limite inferior da instabilidade que uma implantação veria entre sessões.

\section{Conclusão}

Perguntamos quais camadas de um pipeline RAG produzem defeitos que uma suíte de
testes representativa deixa de detectar, quanto da resposta depende do oráculo e
quanto o rigor custa. Sobre 3.000 invocações de um assistente de código aberto
do domínio de engenharia de software, mutado por 18 operadores cobrindo corpus,
segmentação, indexação, recuperação e prompt, a suíte matou 17 dos 18 mutantes
--- e detectou o mediano por 2,5 dos seus 20 casos avaliáveis. Adequação perto
de um e detecção tão fina não são contradição: são o que um escore de mutação
esconde quando é relatado sem a taxa de morte ao lado.

O teste de mutação transfere para esses sistemas, mas não por analogia. Os
operadores precisam agir sobre o corpus, a segmentação, o índice e a
recuperação, onde as falhas documentadas estão; cada um precisa carregar uma
checagem de que age de fato; e o escore precisa ser calculado sobre os casos que
o sistema não-mutado responde de forma estável, porque não há outro terreno
sobre o qual uma morte possa ser atribuída. O que a transposição então revela é
que o número depende de a quem se pergunta: o oráculo mais difundido na prática
não consegue ver um número errado numa sentença correta, e a métrica que a área
trata como padrão não consegue ver uma invenção sobre um tópico que o contexto
não contém. Não se trata de critério de desempate entre instrumentos, e sim de
22 pontos de diferença no número de manchete, com o instrumento mais barato
estritamente dominado por aquele que não custa nada para executar.

\section*{Disponibilidade dos dados}

O pacote de replicação --- sistema sob teste, corpus com hashes de conteúdo e
script de reconstrução, suíte de testes, operadores com suas checagens de
fidelidade, pares de calibração e sua revisão, log bruto de execução, vereditos
e scripts de análise --- será depositado com DOI sob licença Apache-2.0.

\bibliographystyle{plainnat}
\bibliography{refes}

\end{document}
